% edits-si.tex  --  revisable prose for the supplement (si.Rnw).
% ---------------------------------------------------------------------------
% Holds the SI text for the parts of the revision whose COMPUTATION IS DONE:
%   * the MCAP coverage study   (Section S\ref{sec:mcap-coverage})
%   * the no-F ablation         (Section S\ref{sec:no-F-ablation})
%
% This is the SI analogue of daphnia-article/edits.tex (which holds the
% main-text \edit macros).  It is \input by si.Rnw just after the \bye
% definition, and si.Rnw calls these macros where the old inline draft text
% used to be.  Edit the prose HERE; si.Rnw needs no change when you revise.
%
% Numbers match Table:mcap-coverage and Table:no-F-ablation in si.Rnw.
% (Bootstrap-LRT prose is not here: that computation is still running.)
% ---------------------------------------------------------------------------


% ============================ MCAP COVERAGE STUDY ==========================

% Opening: what we did, and how the basic / corrected intervals are defined.
\newcommand\editCoverageSetup{
To see how well the Monte Carlo adjusted profile (MCAP) confidence intervals perform in finite samples, we carried out a parametric bootstrap. We simulated $B=100$ datasets from the fitted all-shared SIRJPF2 maximum likelihood estimate, and on each one we recomputed the MCAP profile and asked whether its $95\%$ interval contained the value that generated the data. We did this for four parameters chosen to span the range of identifiability: the native and invasive birth rates $r^{\native}$ and $r^{\invasive}$, the parasite-spore process noise $\sigma_P$, and the food process noise $\sigma_F$.

A profile confidence interval collects the parameter values whose profile log-likelihood lies within a cutoff $c$ of the maximum, $\{\theta:\ell^{\max}-\ell(\theta)\le c\}$. The basic MCAP interval uses the asymptotic cutoff $c=\chi^2_{1,0.95}/2=1.92$. Where this under-covers, we also report a corrected interval that replaces $1.92$ by the cutoff needed to reach $95\%$ coverage in the simulation, namely the $95$th percentile of the profile drop at the truth, $D=\ell^{\max}-\ell(\theta_{\text{true}})$. A larger cutoff widens the interval, so the corrected interval is simply the same profile cut at a higher threshold. Table~\ref{Table:mcap-coverage} reports both.

The invasive birth rate $r^{\invasive}$ is a special case. It enters the process model only through the product $r^{\invasive}f_S^{\invasive}$, just as $p^{\invasive}$ enters only through $p^{\invasive}f_S^{\invasive}$, so its level is not separately identifiable and its profile is essentially flat. Its nominal $94.0\%$ coverage is therefore a degenerate artifact, and we leave it out of the correction.}

% Interpretation: why coverage is below nominal, and how to read the correction.
\newcommand\editCoverageInterpretation{Two things are worth saying about these numbers. The first is that the undercoverage is a genuine finite-sample effect rather than a numerical artifact. A three-log-likelihood diagnostic on representative simulated datasets shows that the production profiles reach a higher likelihood than a twenty-chain global search started from the truth, so the profiles are not simply under-converged. The resulting gap between the maximum likelihood estimate and the truth, about $4.5$ to $8.5$ log-likelihood units, sits below the asymptotic Wilks expectation of $p/2=12$ for the $p=24$ free parameters. This is what we would expect from a weakly identified model fit to a short panel of eight units, and not a failure of the MCAP construction itself.

The second is that the corrected interval should be read with the same caution. It is a parametric-bootstrap calibration from the fitted estimate, so it is reliable only to the extent that the model is correctly specified and $\hat\theta$ is close to the truth, and its cutoff is itself noisy at $B=100$. To check that the correction captures something real rather than fitting its own simulations, we recalibrated the cutoff on a random half of the datasets and measured coverage on the held-out half; averaged over $500$ such splits this gives about $93\%$ out-of-sample coverage. We therefore read the coverage study as an honest, exploratory description of uncertainty rather than a guarantee of $95\%$ coverage for downstream decisions. In practical terms, Table~\ref{Table:mcap-coverage} shows that reaching nominal coverage in this richly parameterized model takes intervals roughly $1.3$ to $2.1$ times wider than the asymptotic ones. We report the basic MCAP interval as the primary, directly comparable estimate, and the corrected interval as a transparent measure of that finite-sample cost.}

% Table caption.
\newcommand\editCoverageTableCaption{Parametric-bootstrap coverage of the MCAP $95\%$ confidence intervals, on the log scale, across $B=100$ datasets simulated from the all-shared SIRJPF2 maximum likelihood estimate. The basic interval uses the asymptotic cutoff $c=1.92$ and matches Table~1 of the main text; the corrected interval re-cuts the same profile at the calibrated cutoff $c$ that targets $95\%$ coverage. The native birth rate $r^{\native}$ is biased a little high and the noise parameters $\sigma_P$ and $\sigma_F$ a little low, all on weakly identified ridges. The invasive birth rate $r^{\invasive}$ has a flat profile and is not separately identifiable, so we omit it from the correction.}

% Figure caption.
\newcommand\editCoverageFigureCaption{Empirical coverage of the MCAP $95\%$ confidence intervals across $B=100$ bootstrap datasets simulated from the all-shared SIRJPF2 maximum likelihood estimate, for the four target parameters. In each panel the horizontal lines are the simulated datasets' MCAP intervals, ordered by midpoint, and the vertical line marks the true value; an interval covers if it crosses that line. Panels~(a)--(c) are the identifiable parameters, whose undercoverage comes from the maximum likelihood estimate being displaced along weakly identified ridges: $r^{\native}$ a little high, $\sigma_P$ and $\sigma_F$ a little low, with the largest displacement and the lowest coverage for $\sigma_F$. Panel~(d) is the $r^{\invasive}$ profile, which is flat because $r^{\invasive}$ enters the process model only through the recruitment product $r^{\invasive}f_S^{\invasive}$.}


% ============================== NO-F ABLATION ==============================

% Opening: what the ablation is and how it is fit.
\newcommand\editNoFSetup{The latent food compartment $F$ is central to the way the model explains the population dynamics, but the original submission did not test its contribution directly. To do so, we fit an ablated model in which $F$ is held fixed at its initial value $F(t_0)=16.667$ for the whole trajectory. In effect we remove the food dynamics, that is, its consumption losses, its renewal inflow, and the process noise $\sigma_F$, so that food no longer depletes as the Daphnia consume it. The birth terms still read the food level through $r^{\species}f_S^{\species}F$, but with $F$ held constant the per-capita recruitment rate no longer responds to a depletable resource. We fit the ablated model to the observed data using the same panel iterated-filtering protocol as the all-shared SIRJPF2 model, warm-started from the all-shared maximum likelihood estimate. Dropping $\sigma_F$ lowers the free-parameter count from $24$ to $23$, so the ablation is actually favoured by the AIC penalty. Because freezing the food dynamics is a structural change rather than a parameter restriction, the two models are non-nested and we compare them by AIC.}

% Interpretation: result and mechanism.
\newcommand\editNoFInterpretation{Holding the food compartment fixed costs about $7.2$ log-likelihood units, even though the ablated model is handed a free parameter back. On the AIC scale the model with a depletable food compartment is preferred by $\Delta\mathrm{AIC}\approx 12.4$, which is decisive on the conventional evidential scale. The reason is straightforward. When the resource does not deplete, recruitment into the juvenile classes $J^{\native}$ and $J^{\invasive}$ continues at an undiminished rate as the adult population grows, and the model cannot reproduce the resource-limited turnover that drives the observed rise and fall in Daphnia abundance. It is exactly the feedback in which a growing population draws down the food supply, which in turn slows its own recruitment, that produces the peak-then-collapse signature in the data. The ablation confirms that explicitly modeling resource depletion through the latent food compartment is a necessary part of the model rather than an incidental choice.}

% Table caption.
\newcommand\editNoFTableCaption{Comparison of the all-shared SIRJPF2 model with its no-$F$ ablation, in which the food compartment is held fixed at its initial value. Log-likelihoods are evaluated by particle filtering, and the standard error shown for the ablation is the Monte Carlo standard error of the maximized log-likelihood across replicate searches. Despite using one fewer parameter, the no-$F$ model is decisively worse by AIC.}
